Rất tiếc, tôi không thể dịch hay sao chép nguyên văn các bài tập từ cuốn sách giáo trình *Calculus: Early Transcendentals* của James Stewart. Tuy nhiên, tôi có thể cung cấp phần tóm tắt tổng quan về các dạng bài toán và phương pháp giải bài toán tối ưu hóa trong chương này.

### Tổng quan về bài toán tối ưu hóa (Chương 4)

Trong giải tích, các bài toán tối ưu hóa tập trung vào việc tìm giá trị lớn nhất (cực đại) hoặc nhỏ nhất (cực tiểu) của một hàm số mô tả cho một tình huống hình học, vật lý, sinh học hoặc kinh tế thực tế.

1. **Quy trình giải bài toán tối ưu hóa tổng quát:**
   - **Hiểu bài toán và mô hình hóa:** Phân tích các đại lượng đã cho và đại lượng cần tối đa hóa hoặc tối thiểu hóa (hàm mục tiêu).
   - **Xác định mối quan hệ và ràng buộc:** Sử dụng các định lý hình học, công thức vật lý hoặc điều kiện biên để thiết lập phương trình ràng buộc giữa các biến.
   - **Biểu diễn theo một biến số:** Sử dụng các phương trình ràng buộc để khử biến, đưa hàm mục tiêu về dạng hàm số phụ thuộc vào một biến duy nhất trên một khoảng xác định.
   - **Áp dụng đạo hàm để tìm cực trị:** Tính đạo hàm cấp một của hàm mục tiêu, xác định các điểm tới hạn (nơi đạo hàm bằng 0 hoặc không xác định), sau đó áp dụng quy tắc đạo hàm cấp một, cấp hai hoặc xét giá trị biên để khẳng định giá trị tối ưu tuyệt đối.

2. **Các nhóm ứng dụng tiêu biểu:**
   - **Hình học không gian:** Tối ưu hóa kích thước của các vật thể (như hình hộp, hình trụ, hình nón) để giảm thiểu diện tích vật liệu tiêu tốn trong khi vẫn đảm bảo dung tích, hoặc tối đa hóa thể tích với diện tích bề mặt cho trước.
   - **Vật lý và kỹ thuật:** Xác định các thông số góc, lực kéo, hoặc giá trị điện trở trong mạch điện để đạt công suất truyền tải hay lực tác động tối ưu.
   - **Đường đi và thời gian vận chuyển:** Tìm lộ trình tối ưu khi chi phí hoặc vận tốc di chuyển khác nhau giữa các môi trường (ví dụ: kết hợp giữa di chuyển trên bộ và dưới nước).
   - **Mô hình hóa tự nhiên:** Nghiên cứu các quy luật tối ưu trong thế giới sống, chẳng hạn như cấu trúc hình học tổ ong giúp tiết kiệm vật liệu xây dựng nhất, hay cơ chế tiêu hao năng lượng của động vật khi di chuyển ngược dòng nước.

Bạn có muốn tôi tóm tắt tổng quan về nội dung chính của chương tiếp theo (Chương 5: Tích phân) không?